\documentclass[a4paper]{article}

% Basic packages
\usepackage[fontset=fandol]{ctex}
\usepackage{amsmath, amssymb}
\usepackage{graphicx}
\usepackage[margin=2.5cm]{geometry}
\usepackage[most]{tcolorbox}
\usepackage{etoolbox}
\usepackage{listings}
\usepackage{hyperref}
\usepackage{booktabs}
\usepackage{subcaption}
\usepackage{float}
\usepackage{tikz}
\IfFileExists{pgfplots.sty}{
    \usepackage{pgfplots}
    \pgfplotsset{compat=1.18}
}{}

% Highlight boxes
\newtcolorbox{knowledgebox}[1]{
    enhanced, colback=blue!5!white, colframe=blue!75!black, colbacktitle=blue!75!black,
    coltitle=white, fonttitle=\bfseries, title=#1, attach boxed title to top left={yshift=-2mm, xshift=2mm},
    boxrule=1pt, sharp corners
}

\newtcolorbox{importantbox}[1]{
    enhanced, colback=yellow!10!white, colframe=yellow!80!black, colbacktitle=yellow!80!black,
    coltitle=black, fonttitle=\bfseries, title=#1, sharp corners
}

\newtcolorbox{warningbox}[1]{
    enhanced, colback=red!5!white, colframe=red!75!black, colbacktitle=red!75!black,
    coltitle=white, fonttitle=\bfseries, title=#1, sharp corners
}

% Listings
\lstset{
    language=Python,
    basicstyle=\ttfamily\small,
    keywordstyle=\color{blue},
    stringstyle=\color{red!60!black},
    commentstyle=\color{green!60!black},
    breaklines=true,
    frame=single,
    numbers=left,
    numberstyle=\tiny\color{gray},
    captionpos=b,
    extendedchars=false
}

% Front-page metadata
\newcommand{\notetitle}{Lecture 9: Case Studies in Locomotion}
\newcommand{\noteauthors}{Codex}
\newcommand{\notedate}{2026-04-16}
\newcommand{\videochannel}{Toru}
\newcommand{\videopublishdate}{2024-11-04}
\newcommand{\videoduration}{1:26:05}
\newcommand{\videourl}{https://www.youtube.com/watch?v=98uo2OuXDYc}
\newcommand{\videocoverpath}{cover.jpg}

\begin{document}

\begin{titlepage}
\centering
{\Large 课程笔记\par}
\vspace{1.2cm}
{\huge\bfseries \notetitle\par}
\vspace{0.8cm}
{\large \noteauthors\par}
\vspace{0.3cm}
{\large \notedate\par}
\vspace{1.2cm}

\ifdefempty{\videocoverpath}{
{\small 请在生成笔记时填入视频封面图的本地路径。\par}
}{
\includegraphics[width=0.82\textwidth,height=0.45\textheight,keepaspectratio]{\videocoverpath}\par
}

\vfill
\begin{tcolorbox}[width=0.9\textwidth, colback=black!2!white, colframe=black!60, sharp corners]
\textbf{视频作者/频道}：\videochannel\par
\textbf{发布日期}：\videopublishdate\par
\textbf{视频时长}：\videoduration\par
\textbf{视频链接}：
\ifdefempty{\videourl}{
未填写
}{
\href{\videourl}{\nolinkurl{\videourl}}
}
\end{tcolorbox}
\end{titlepage}

\tableofcontents
\newpage

\section{四足 locomotion：从 pattern generator 到强化学习}
\subsection{问题从哪里来}
四足机器人（quadruped）看起来比 humanoid 更容易，但真正能在现实中走稳、走快、跨越复杂地形，依然非常难。课程先用自然界里的狗、马和山地地形作为直观背景，再回到机器人：腿式系统需要同时面对平衡、能量、接触切换和地形变化。

\subsection{早期方法：Pattern Generator 与 Raibert Controller}
早期腿式机器人控制强调 \textbf{pattern generator（节律发生器）} 和固定步态。核心思想是用预编程的中央模式发生器（central pattern generator, CPG）或 Raibert controller 产生重复、稳定、可解释的步态。它们的优点很明确：实现简单、实时性强、在结构化环境里够用。

但这类方法也有明显局限：它们依赖人工设计的步态和环境假设，一旦地形、摩擦、负载或速度目标变化，就需要重新调参。课程在这里给出的判断非常直接：\textbf{固定步态能走，但不够“会走”}。

\begin{figure}[H]
\centering
\includegraphics[width=0.92\textwidth,page=14]{slides.pdf}
\caption{中央模式发生器（CPG）是早期腿式 locomotion 的代表思路之一。}
\end{figure}

\subsection{RL 的转折}
现代四足 locomotion 的核心转折是：让机器人在模拟器里用 \textbf{reinforcement learning (RL)} 自己学步态。课程里提到的关键奖励通常围绕两件事：\textbf{最小化工作量（work）} 和 \textbf{最小化地面冲击（ground impact）}。这和生物运动学里的能量规律是一致的，因此会自然涌现 walking、trotting、galloping 等不同 gait。

\begin{importantbox}{为什么 RL 在 locomotion 里好用}
locomotion 的一个优点是奖励相对容易写：走得稳、走得快、摔得少、能耗低。这使 RL 比 manipulation 更容易先跑起来，也更容易在 simulation 里大规模并行训练。
\end{importantbox}

\subsection{本章小结}
四足 locomotion 的主线是从固定步态走向可学习步态。pattern generator 提供了早期可用的工程解，但 RL 才让机器人能根据速度、地形和扰动自适应地组织 gait。

\section{从仿真到现实：RMA 与速度自适应}
\subsection{Rapid Motor Adaptation}
课程把 \textbf{Rapid Motor Adaptation (RMA)} 作为从 simulation 到 real world 的代表方法。其关键不是只训练一个 policy，而是把系统拆成两层：

\begin{itemize}
  \item \textbf{Base policy}：在模拟器中学会基本的行走控制。
  \item \textbf{Adaptation module}：根据观察到的历史状态，在线估计隐含的环境参数，例如摩擦、地形高度、载荷和地面柔软度。
\end{itemize}

在训练阶段，RMA 会利用 privileged terrain information 监督地形编码器；在部署阶段，这些 privileged 信息不可用，于是 adaptation module 只能根据过去的 observation-action history 推断 extrinsics。课程强调这一步实际上是一个 \textbf{system identification（系统辨识）} 风格的问题，只不过它是 learned 的、在线的。

\begin{figure}[H]
\centering
\includegraphics[width=0.92\textwidth,page=27]{slides.pdf}
\caption{RMA 的核心直觉：用历史观测在线估计不可见的环境外参（extrinsics）。}
\end{figure}

\subsection{速度不同，步态就不同}
另一条线索是 \textbf{speed-adaptive gaits（速度自适应步态）}。课程展示了一个非常重要的现象：当训练目标从单一速度变成多个速度时，机器人会自发学出不同步态。低速更像 walk，中速更像 trot，高速则可能出现更明显的 bounce/gallop。这里不是人工“规定”步态，而是奖励函数与动力学共同诱发了步态分化。

讲义还强调了一个生物学上的对应关系：最小化能量消耗会和真实动物的 gait pattern 对齐。也就是说，机器人并不是在“模仿动物动作表面”，而是在重现某些深层的优化规律。

\subsection{本章小结}
RMA 证明了一个实用结论：如果环境外参可以从历史中估计，那么很多原本只能在仿真里训练的策略就能被搬到现实。速度自适应则进一步说明，步态并不是固定模板，而是可以由任务目标和能量约束自然生成的。

\section{视觉与本体感觉：腿式机器人的导航}
\subsection{为什么还需要 vision}
课程从 locomotion 走向 navigation 时，重点是补上 perception。仅靠 proprioception，机器人可以在一定程度上维持平衡，但遇到 stairs、curbs、stepping stones、slopes、glass walls 这样的复杂场景时，单靠内部状态并不够。于是需要 \textbf{egocentric vision（第一人称视觉）}，尤其是 depth 信息。

这里有一个很重要的现实判断：直接建 terrain map 在理论上很自然，但在现实中很脆弱，因为机器人自身和相机都在抖，累积噪声会让地图逐步失真。课程因此提出一个更直接的问题：\textbf{我们真的需要 terrain map 吗？}

\begin{figure}[H]
\centering
\includegraphics[width=0.92\textwidth,page=95]{slides.pdf}
\caption{从视觉直接到控制，是这部分内容的转折点。}
\end{figure}

\subsection{地图化与 map-free 方法}
讲义先讲传统的 map-based perceptive locomotion：利用视觉构造地形图，再把地图输入低层 locomotion policy。随后给出的替代路线是 \textbf{map-free, gait-free}：不显式维护精细 terrain map，而是直接从视觉和历史状态中学习控制。课程展示的结果说明，在某些挑战地形上，这种方法反而比 noisy map 更稳。

\begin{knowledgebox}{这部分的中心结论}
对于腿式机器人，感知的目标不是“重建一个完美世界模型”，而是“尽快提取对下一步足部放置有用的信息”。因此，地图并非唯一答案，甚至在某些场景里会是负担。
\end{knowledgebox}

\subsection{本章小结}
vision 在腿式 locomotion 里的作用，不是让机器人更像传统导航系统，而是让控制系统更像一个可以实时理解地形的运动体。课程最后的经验结果表明，地图未必总是收益，直接从视觉到控制有时更有效。

\section{双足与 humanoid：ZMP、MPC 与学习式控制}
\subsection{为什么 humanoid 更难}
与四足相比，双足机器人更接近人类，但稳定性也更差。课程先用 ASIMO、ATLAS 等例子说明传统 humanoid 控制如何依赖 \textbf{zero-moment point (ZMP)} 和 \textbf{model-predictive control (MPC)}。这些方法在结构化环境里非常强，因为它们依赖明确的动力学模型和可预测的接触。

问题在于，一旦环境变得非结构化，或者模型不够准确，纯模型法就会迅速吃力。换句话说，humanoid 的难点不只是“做出一个步态”，而是“在现实里持续地、鲁棒地走”。

\subsection{端到端 RL 控制}
课程随后转向 Radosavovic et al. 的工作：用 \textbf{causal transformer} 和 \textbf{reinforcement learning} 训练一个端到端 humanoid controller。输入是历史 proprioception（关节角、速度等），输出是关节残差控制。多环境、不同摩擦、不同地形的 domain randomization 让策略学到在 context 下自适应的行为。

这部分最值得注意的现象包括：

\begin{itemize}
  \item arm swing 并没有显式写进奖励，但会自发涌现；
  \item 在坡地上，步长会变小、步频会提高；
  \item 当机器人遇到台阶或扰动时，会出现 foot recovery 行为。
\end{itemize}

\begin{figure}[H]
\centering
\includegraphics[width=0.92\textwidth,page=117]{slides.pdf}
\caption{端到端 humanoid control 的典型形态：直接学控制，而不是手写 ZMP/MPC 细节。}
\end{figure}

\subsection{本章小结}
humanoid locomotion 的重要变化是：从模型优先转向学习优先。ZMP 和 MPC 仍然有价值，但在现实复杂性面前，学习式控制提供了更强的适应性，尤其是在足够大的仿真训练和合理的 domain randomization 之后。

\section{从视频到序列：next-token prediction 与 challenging terrains}
\subsection{Humanoid Locomotion as Next Token Prediction}
课程后半段把 humanoid locomotion 重新表述成一个序列建模问题：像语言模型预测下一个 token 一样，预测下一个动作片段或控制信号。这是一个非常强的类比，因为它把 locomotion 从“在线控制”改写成“序列建模 + in-context adaptation”。

这类方法通常先用 human walking videos 和 mocap 数据做 \textbf{kinematic retargeting（运动学重定向）}，再在大规模数据上训练 transformer。课程强调了一个很现实的取舍：数据可以 noisy，但只要 scale 足够大，模型仍然能学到有用的 gait priors。

\begin{figure}[H]
\centering
\includegraphics[width=0.92\textwidth,page=136]{slides.pdf}
\caption{训练中引入 missing data，是把 locomotion 写成序列建模问题的关键手段之一。}
\end{figure}

\subsection{预训练再微调}
在 challenging terrain 上，课程展示了另一条混合路线：先用 sequence modeling / imitation 做预训练，让机器人先“会走”；再用 RL fine-tuning 去适应更难的 terrain。直觉很简单：一个已经能走平地的机器人，比一个完全不会走的机器人更容易学会上坡、跨越障碍和应对脚步失误。

这类方法的优势和代价都很清楚：

\begin{itemize}
  \item 优势是 sample efficiency 高，训练更稳定；
  \item 代价是它只能学习数据里存在的技能，真正全新的技能仍然更依赖 RL。
\end{itemize}

\begin{importantbox}{课程最后的态度}
对 locomotion 而言，imitation learning 和 RL 不是互斥的。更实用的范式往往是：先用大规模视频和序列建模把策略带到一个好初值，再让 RL 处理数据里没有覆盖到的鲁棒性和新技能。
\end{importantbox}

\subsection{本章小结}
这部分把 locomotion 的前沿总结得很清楚：一条线是更强的感知与控制闭环，另一条线是把动作当作长序列来建模。二者最终都在追求同一个目标，即让机器人在真实世界里以更少样本、更高稳定性学会走路。

\section{总结与延伸}
这讲的核心结论可以概括为三层。第一层，\textbf{quadruped locomotion} 证明 RL 可以学出自然步态，且能把 energy-efficient gait 与生物规律联系起来。第二层，\textbf{RMA} 和 vision-proprioception coupling 说明，现实部署的关键是处理不可见外参和感知噪声，而不是只把 simulator 的 policy 原封不动搬过去。第三层，\textbf{humanoid locomotion} 显示，未来会越来越像序列建模问题：先在大规模 noisy data 上学出先验，再用 RL 处理不可预见的环境与任务。

如果把这一讲和前后的课程放在一起看，它其实在反复强调同一个模式：\textbf{机器人学习最有效的路径，不是去掉结构，而是用学习把结构用起来。} locomotion 里这个结构是动力学、接触和能量；视觉里是人体先验和可供性；navigation 里则是语义地图和分层控制。

从教材角度看，这一讲值得记住的不是某一个单独模型，而是三个原则：一，simulation 是训练大型 policy 的基础设施；二，现实部署需要适应机制；三，越接近真实世界，越需要把感知、控制和任务结构放在同一个框架里思考。

\end{document}
